\chapter{Methods}

\section{First task}
\subsection{Model}
\todo{explain problem and model}

\subsection{Operator Splitting}

Consider a typical partial differential equation (PDE) of the form
$$ \partial_t \bm{u} = (A + B + C)\bm{u}. $$
Here, $\bm{u}$ denotes the fluid velocity at a given location. The terms
$A$, $B$, and $C$ are called \textit{operators}, which represent 
individual processes within the PDE corresponding to different physical 
effects. As an example, consider the basic advection-diffusion equation
$$\partial_t T = \Delta T - \bm{u} \cdot \nabla T.$$
Defining the operators $A := \Delta$ and $B := \bm{u} \cdot \nabla$, the 
equation can be written as
$$\partial_t T  = (A - B) T.$$
Here, $A$ represents the Laplace operator (diffusion), while $B$ 
represents advection by the fluid velocity field $\bm{u}$.

For most PDEs, analytical solutions do not exist, and one must instead 
use numerical methods to approximate their solutions. However, these 
methods can be computationally expensive. \textit{Operator splitting} 
reduces this computational cost by partitioning the problem into simpler
subproblems, where each operator is treated individually. To continue 
with our previous example, one solves the following subproblems 
sequentially: 
\begin{align}
\partial_t \bm{u} &= A \bm{u} \\
\partial_t \bm{u} &= B \bm{u}
\end{align}


When the first equation is solved numerically, the resulting state is then
used as the initial condition for the second subproblem. This procedure is
repeated at each timestep. However, this introduces an approximation 
error. The intermediate state does not exactly correspond to the solution
that would be obtained by analytically solving the combined system.
Different operator splitting methods lead to different magnitudes of
error. In general, the error remains controlled if the timestep $\Delta t$
is sufficiently small.

\subsubsection{First Order Operator Splitting}

First-order operator splitting solves each operator sequentially over a 
timestep $\Delta t$. To illustrate this, consider the equation 
\begin{align}
    \partial_t  \bm{u} &= (A+B)\bm{u} \\
\end{align}
Instead of solving the full equation directly, the problem is approximated
by solving the following subproblems in sequence:
\begin{align}
    \partial_t \bm{u} &= A\bm{u}, \\
    \partial_t \bm{u} &= B\bm{u}. \\
\end{align}

For our problem, these operators arise from the advection-diffusion 
equation, as will be discussed later. Note that applying operator 
splitting results in two separate \textbf{linear} subproblems. This allows
us to obtain simple solutions for the resulting split operator problems,
namely
\begin{align}
    \bm u_A(t) &= e^{(t-t_0) A}\bm u_0 \\
    \bm u_B(t) &= e^{(t-t_0) B}\bm u_0 
\end{align}
In practice, however, we only advance the solution over the next discrete
time interval $[t_k, t_{k+1}]=[t_k, t_k + \Delta t]$, rather than solving
continuously over the entire time domain. Therefore, the solutions are
rewritten for a single timestep as
\begin{align}
    \textbf u_{k+1, A} &= e^{(t_{k+1} - t_k) A}\textbf u_{k, A} \\
    \textbf u_{k+1, B} &= e^{(t_{k+1} - t_k) B}\textbf u_{k, B} 
\end{align}
Using $\Delta t = t_{k+1}-t_k$, we obtain the timestep formulation

\begin{align}
    \textbf u_{k+1, A} &= e^{\Delta t A}\textbf u_{k, A} \\
    \textbf u_{k+1, B} &= e^{\Delta t B}\textbf u_{k, B},
\end{align}
where an equidistant discretization in time is assumed. 

For first-order operator splitting, the following algorithm is applied to
the abstract problem $\partial_t \bm{u}= (A+B)\bm{u}$. The equation is
split into two separate operator problems:
\begin{enumerate}
    \item{Set $\bm{u}_{k,A}:= \bm{u}_k$.}
    \item{Solve $\bm{u}_{k+1, A} = e^{\Delta t A}\bm u_{k, A}$ on the
            interval $[t_{k}, t_k+\Delta t].$}
    \item{Use $\bm{u}_{k+1, A}$.}    
    \item{Consider $\bm{u}_{k+1, A}$ as the initial condition for the 
            second operator problem by setting
            $\bm{u}_{k,B}:=\bm{u}_{k+1,A}.$}
    \item{Solve $\bm{u}_{k+1, B} = e^{\Delta t B}\bm{u}_{k, B}
            =e^{\Delta t B}\bm{u}_{k+1, A}$.}
    \item {Set $\bm{u}_{k+1} := \bm{u}_{k+1,B}$.}
\end{enumerate}
Written in a compact form, the procedure becomes
$$\bm u_{k+1}= e^{\Delta t B}e^{\Delta t A}u_k$$
As previously mentioned, this algorithm relies on the assumption that for
small timesteps $\Delta t$, only small difference remains between 
$u_{k+1,A}$ and $u_{k+1,B}$ when the operators are applied subsequently.
This algorithm is simple and easy to implement. 
\todo{include some computational efficiency stuff, maybe even related to
our testing?}
However, the global truncation error (GTE) is of order
$\mathcal{O}(\Delta t)$, \todo{include source} which motivates the
development of second-order splitting methods.

\subsubsection{Second Order Operator Splitting (Strang-Splitting)}
To reduce the GTE of the operator splitting scheme to order
$\mathcal{O}(\Delta t^2)$, Strang splitting can be applied. The solution
$\bm{u}_{k+1}$ at the next timestep $t_{k+1}=t_k +\Delta t$ is
approximated by
\begin{equation}
    \bm{u}_{k+1}=e^{\frac{\Delta t}{2}A} e^{\Delta t B}
        e^{\frac{\Delta t}{2}A} \bm{u}_k
\end{equation}

\subsection{The Velocity Equation}
We have seen operator splitting for a general problem, and now apply the
method to the context of the velocity equation from the Navier-Stokes
equations
$$ \partial_t \bm{u} + \bm{u} \cdot \nabla \bm u + \frac{1}{\rho}\nabla p
= \nu \Delta \bm{u} + \beta T\bm{e}_z $$
This equation can be split into three different operators:
\begin{enumerate}
    \item{Advection: $\partial_z \bm{u} = \bm{u}\cdot \nabla \bm{u}$}
    \item {Diffusion + forcing: $\partial_t \bm {u} = \nu \Delta \bm{u}
            + \beta T_{\bm{e}_z}$}
    \item {Projection step (Incompressibility): $\bm{u}^{k+1}= \bm{u}^*-
            \Delta \bm{p}, \quad \text{with } \nabla \cdot \bm{u}^{k+1}=0$
             \\ The pressure acts like a Lagrange multiplier, enforcing
             the incompressibility constraint, $\nabla \bm{u}=0.$}
    
\end{enumerate}

To compute the numerical solution for the velocity field, the problem is
solved in two stages.
\begin{enumerate}
    \item{\textbf{Step 1: Compute intermediate velocity} \\ Ignoring the
        incompressibility constraint, we solve $\partial_t \bm{u}+\bm{u}
        \cdot \nabla \bm{u} = \nu \cdot \Delta \bm{u} + f$. This yields an
        intermediate velocity $\bm u^*$ that generally does not satisfy
        the divergence-free condition, $\nabla \cdot \bm u^* \neq 0$}
    
    \item{\textbf{Step 2: Projection step} \\ The intermediate velocity is
    then corrected to enforce incompressibility: $$ \bm{u}^{k+1} = 
    \bm{u}^*-\nabla \bm{p}.$$ Using the \textit{Helmholtz decomposition},
    $$\bm{u}^* = \bm{u}^\perp + \nabla \phi,$$ where $$\bm{u}^{\perp}$$ is
    divergence free. Taking the divergence yields $$\nabla \cdot \bm{u}^* 
    = \Delta \phi.$$ Imposing incompressibility on $\bm{u}^{k+1}$, 
    $$\nabla \cdot \bm{u}^{k+1} = 0,$$ leads to the Poisson equation for
    the pressure: $$\Delta p = \nabla \cdot \bm{u}^*.$$}
\end{enumerate}


Here, we briefly summarize the developments so far. Readers who feel
confident with the material presented above may skip the following
paragraphs. Nevertheless, we provide a short overview of the underlying
problem to make the discussion in the later sections easier to follow.

We begin by giving context to the physical problem addressed by the
Navier--Stokes equations.

The Navier-Stokes equations are partial differential equations that are
generally difficult to solve analytically. Consequently, numerical methods
are necessary to compute the velocity on our domain.

The physical problem described by the Navier-Stokes can be interpreted as
follows. Certain properties of the fluid must be satisfied, such as how
the volume changes under pressure, how the flow evolves with time, how the
motion at one point is influenced by surrounding fluid, etc. These 
properties are represented by different operators that describe the 
evolution of the flow.
Our main interest, therefore, is the velocity of the fluid at each grid
point in our domain at each timestep. The temporal development of the 
velocity at a given location is determined by spatial characteristics at
that location, such as the first and second spatial derivatives of the 
velocity relative to the surrounding fluid. In other words, the 
\textit{temporal} evolution is driven by \textit{spatial} derivatives of
the velocity itself. 
Advection, for example, describes the transport of fluid along the 
direction of flow. Diffusion, on the other hand, models the smoothing of
velocity differences within the domain. Therefore, in the absence of any
external force or energy, input would cause unevenly distributed velocity
behavior, and we would observe an eventual state of rest.
These mechanisms depend only on characteristics of the velocity field 
$\textbf u$. However, additional terms such as buoyancy forcing 
$\beta T_z$ and pressure gradient $\frac{1}{\rho}\nabla p$ introduce
additional complexity. In the following sections, we neglect additional
forcing terms (which could represent, for example, heating to induce
convection) and focus primarily on the effects of the pressure term.

Consider a generalized incompressible fluid, characterized by $\nabla 
\cdot \textbf u$ = 0. By definition, an incompressible fluid does not
change its volume when forces or pressure are applied. Consequently, no
volume should be lost or gained during the flow. Pressure, which is
coupled to the velocity field, therefore becomes an additional dynamical
variable that must be solved for.

If the governing equation only depends on the velocity field $\textbf{u}$,
we could simply discretize the equation and solve numerically for 
$\textbf{u}$. But with the presence of the pressure term $p$, this makes
the problem more complicated. To address this, we break the numerical
process into two stages. First, the coupling between velocity $\textbf u$
and pressure $p$ is temporarily ignored and the incompressibility
condition $\nabla \textbf u = 0$ is not enforced. This step yields an
intermediate velocity $\textbf u^*$ for the velocity for the next
timestep. Secondly, the pressure is introduced to enforce the
incompressibility constraint. This can be done efficiently by applying the
\textit{Helmholtz decomposition} to the velocity field. As a result, the
pressure can be determined by solving the Poisson equation, which ensures
the corrected velocity field satisfies the divergence-free condition.

We have obtained numerical solutions for an intermediate velocity and
pressure that satisfies the incompressibility constraint. Using these
quantities, we can compute the velocity at the next discrete timestep
while preserving incompressibility. A detailed description of this
procedure is given in Chapter~\ref{Chorins}.

\subsubsection{Operator splitting for the Velocity Equation}

Returning to the operators of the velocity equation, we now apply
operator splitting for all the different terms.
We begin with the advection operator
$$\partial_t \bm{u} + a_y\partial_y \bm{u} + a_z\partial_z\bm{u} = 0$$
Here we define
$
\bm{a} = 
\begin{pmatrix}
    a_y(y,z) \\
    a_z(y,z)
\end{pmatrix}$
and 
$\bm{u} = \begin{pmatrix}
    v \\
    w
\end{pmatrix}$

The operator splitting is first applied in their spatial directions. This
leads to two one-dimensional advection problems:
\begin{enumerate}
    \item{Solve $$\partial_t \bm{u} + a_y\partial_y \bm{u} = 0 $$ Or, more
    explicitly written component wise: 
    $$ \begin{cases} \partial_t v + a_y \partial_y v &= 0 \\
    \partial_t w + a_y \partial_y w &= 0 \end{cases}$$} 
    We do the same for the second spatial dimension:
    \item{Solve $$\partial_t \bm{u} + a_z\partial_z \bm{u} = 0 $$ Again,
    more explicitly written component wise: 
    $$\begin{cases} \partial_t v + a_z \partial_z \bm{u} &= 0\\
    \partial_t w + a_z \partial_z w &= 0 \end{cases}$$} 
    \item{Apply Lie Splitting in time: $\bm{u}^{k+1} = 
        A_z(\Delta t)A_y(\Delta t)\bm{u}^{k}$}
\end{enumerate}

We have now obtained an operator splitting for the (linear) advection
equation. In the next section, we will examine the upward finite
difference method to solve the resulting equations numerically.

\subsection{Upwind Finite Difference Method}
To introduce the upwind finite difference method, we consider
the one-dimensional linear advection equation
$$\frac{\partial_u}{\partial_t}+a\frac{\partial u}{\partial x}=0.$$
This equation describes a wave propagating along the $x$-axis with
velocity $a$. If $a>0$, the wave propagates to the right. Thus, an upwind
direction would correspond to the direction in the left, towards smaller
$x$ values. Conversely, if $a<0$, the wave propagates to the left, and the
upwind direction corresponds to larger values of $x$.

Therefore, we will use backward finite difference methods for $a>0$ and
forward finite difference methods for $a<0$. For the one-dimensional
discretization we define
\[
\begin{aligned}
x_i &:= i \,\Delta x, \quad i = 0,\dots,N \\
t^n &:= n \,\Delta t
\end{aligned}
\]
\todo{also define discrete no of steps for Time? So maybe $N_x$ and
$N_t$?}

\begin{enumerate}
    \item{$a > 0$ (backward): $$\partial_x \bm{u} \approx \frac{\bm{u}_i 
    - \bm{u}_{i-1}}{\Delta x} $$}
    \item{$a < 0$ (forward): $$\partial_x \bm{u} \approx \frac
    {\bm{u}_{i+1} - \bm{u}_{i}}{\Delta x} $$}
\end{enumerate}

Using a forward Euler discretization in time,
$$\frac{\bm{u}_{i+1}^n - \bm{u}_{i}^{n+1}}{\Delta t} 
= -a(\partial_x\bm{u})^n_i $$


\begin{enumerate}
\item $a > 0$ (backward):
\begin{align}
\frac{\bm{u}_{i+1}^n - \bm{u}_{i}^{n+1}}{\Delta t}
  &= -a(\partial_x \bm{u})^n_i \\
\Rightarrow \bm{u}_i^{n+1} - \bm{u}_i^n
  &= -a \frac{\Delta t}{\Delta x}(\bm{u}_i^n - \bm{u}_{i-1}^n) \\
\Rightarrow \bm{u}_i^{n+1}
  &= \left(1 - a\frac{\Delta t}{\Delta x}\right)\bm{u}_i^n
     + a\frac{\Delta t}{\Delta x}\bm{u}_{i-1}^n \\
\Rightarrow \bm{u}_i^{n+1}
  &= A_1 \bm{u}_i^n
\end{align}

\item $a < 0$ (forward):
\begin{align}
\frac{\bm{u}_{i+1}^n - \bm{u}_{i}^{n+1}}{\Delta t}
  &= -a(\partial_x \bm{u})^n_i \\
\Rightarrow \bm{u}_i^{n+1} - \bm{u}_i^n
  &= -a \frac{\Delta t}{\Delta x}(\bm{u}_{i+1}^n - \bm{u}_i^n) \\
\Rightarrow \bm{u}_i^{n+1}
  &= \left(1 - a\frac{\Delta t}{\Delta x}\right)\bm{u}_i^n
     - a\frac{\Delta t}{\Delta x}\bm{u}_{i+1}^n \\
\Rightarrow \bm{u}_i^{n+1}
  &= A_2 \bm{u}_i^n
\end{align}
\end{enumerate}

Thus, we obtain an operator for each upwind case, depending only on the
sign of $a$. To compactly represent the update for all grid points, we
write the scheme in matrix form.

$$ 
A_1 = \begin{pmatrix}
 \left(1-a\frac{\Delta t}{\Delta x}\right) & 0 &  & \dots & 0 \\
 a\frac{\Delta t}{\Delta x} & \ddots &  & & \vdots \\
 0 & \ddots & \ddots & & \vdots \\
 \vdots & & \ddots & \ddots & 0 \\
 0 & \dots & 0 & a\frac{\Delta t}{\Delta x} & \left(
    1-a\frac{\Delta t}{\Delta x}\right) \\
\end{pmatrix} $$

$$ 
A_2 = \begin{pmatrix}
 \left(1-a\frac{\Delta t}{\Delta x}\right) & -a\frac{\Delta t}{\Delta x}
 & 0 & \dots & 0 \\
 0 & \ddots & \ddots & & \vdots \\
 &  & \ddots & \ddots & 0 \\
 \vdots & & & \ddots & -a\frac{\Delta t}{\Delta x}\\
 0 & \dots &  & 0 & \left(1-a\frac{\Delta t}{\Delta x}\right) \\
\end{pmatrix} $$

\subsubsection*{Stability of the Operators}
In general, a matrix operator $A$ is stable if its infinity norm satisfies
$\|A\|_{\infty} \leq 1$. We investigate both operators $A_1$ 
(corresponding to the case $a > 0$ and therefore backward finite
difference method) and $A_2$ (corresponding to the case $a < 0$ and
therefore forward finite difference method).

\begin{align}
    \|A_1\|_{\infty} &= \left|a\frac{\Delta t}{\Delta x}\right| + \left|
        1- a \frac{\Delta t}{\Delta x}\right| \leq 1 \\
    \Rightarrow^{a> 0}  \left|1- a \frac{\Delta t}{\Delta x}\right| &\leq
    1 - a \frac{\Delta t}{\Delta x} \\
    \Rightarrow 1 - a\frac{\Delta t}{\Delta x} &\geq 0 \\
    \Rightarrow a \frac{\Delta t}{\Delta x}&\leq 1 
\end{align}

Repeating this argument for our second operator $A_2$ yields
\begin{align}
    \|A_2\|_{\infty} &= \left|a\frac{\Delta t}{\Delta x}\right| + 
    \left|1- a \frac{\Delta t}{\Delta x}\right| \leq 1 \\
    \Rightarrow^{a> 0}  \left|1- a \frac{\Delta t}{\Delta x}\right| 
    &\leq 1 + a \frac{\Delta t}{\Delta x} \\
    \Rightarrow 1 + a\frac{\Delta t}{\Delta x} &\geq 0 \\
    \Rightarrow - a \frac{\Delta t}{\Delta x}&\leq 1 
\end{align}

Combining both cases yields the general stability condition
$$\left|a\frac{\Delta t}{\Delta x}\right| \leq 1,$$ 
which is the CFL condition for the upwind scheme.

\subsection{Two-Dimensional Upwind Scheme with Operator Splitting}
Now we jointly apply operator splitting with the upwind finite difference
method for the two-dimensional advection equation. We define $u_{ij}^n
:= u(t^n, y_i, z_i)$ and assume an equidistant spatial grid. The spacings
$\Delta y$ and $\Delta z$ denote the distances between neighboring grid
points.

We first solve the advection term in the $y$ direction. For a fixed value
$z_j$, the equation becomes
$$\partial_t \bm{u} + a_y\partial_y \bm{u} = 0.$$
The spatial derivative is approximated using an upwind scheme,
depending on the sign of $a_y$:
\[
(\partial_y\bm{u})_{ij}=
\begin{cases}
    \frac{\bm{u}_{ij}-\bm{u}_{i-1 \,j}}{\Delta y} \quad a_y > 0 \\
     \frac{\bm{u}_{i+1\,j}-\bm{u}_{i\,j}}{\Delta y} \quad a_y < 0 \\
\end{cases}
\]

Using a forward Euler step in time yields
$$u_{ij}^*= u_{ij}^n - \Delta t a_{y,ij}(\partial_y \bm{u})_ij.$$
Stability requires the CFL condition
$$a_y\frac{\Delta t}{\Delta y}\leq 1.$$
We use this result to proceed with the solution of the $z$-coordinate
corresponding to the second operator:
$$\partial_t \bm{u} + a_z\partial_z \bm{u} = 0$$
We define analogously
\[
(\partial_y\bm{u}^*)_{ij}=
\begin{cases}
    \frac{\bm{u}^*_{ij}-\bm{u}^*_{i-1 \,j}}{\Delta y} \quad a_y > 0 \\
     \frac{\bm{u}^*_{i+1\,j}-\bm{u}^*_{i\,j}}{\Delta y} \quad a_y < 0 \\
\end{cases}
\]
Using the result from the first split operator, we now compute
$$u_{ij}^{n+1}= u_{ij}^* - \Delta t a_{z,ij}(\partial_z \bm{u})_ij$$

We ensure stability in this case as well by respecting $\frac{|a_z|
\Delta t}{\Delta z}\leq 1$.

We can now derive a stability condition for all of our discretized points
in space. We first obtain
$$\max |a_y|= \max_{i,j}|a_y(y_i,z_j)|$$
and 
$$\max |a_z|= \max_{i,j}|a_z(y_i,z_j)|.$$
With the help of these variables, we deduce the overall constraint for
$\Delta t$, that is necessary to ensure stability:
$$ \Delta t \leq \min{\left(\frac{\Delta y}{\max|a_y|}, 
\frac{\Delta z}{\max|a_z|}\right)}$$

In summary, we have implemented operator splitting together with the
upwind finite difference method for the two-dimensional linear advection
equation and derived a stability condition that must be satisfied over the
entire computational domain.
\todo{boundary conditions implementation}
\newpage

\section{Incompressible Flow - Chorin's Projection Method}\label{Chorins}

We now consider the Navier-Stokes equations for velocity field
\begin{align}
    \partial_t \textbf{u} + \textbf{u} \cdot \nabla \textbf{u} 
    + \frac{1}{\rho}\nabla p &= \nu \Delta \textbf{u}  \quad \quad 
    \text{with } \beta = 0 \\
    \nabla \cdot \textbf u &= 0  
\end{align}

The boundary conditions are given by
\begin{align}
\textbf u (t,0,z) &= \textbf u (t,1,z)\\
\textbf w (t,y,0) &= \textbf w (t,y,1)=0 \quad \quad \forall t 
\in [0, t_{end}], (y,z)\in \Omega
\end{align}

To solve this differential equation, we employ a \textit{projection
method}.\todo{explain why we use a projection method} Conceptually, this
projection method splits the computation into two steps:
\begin{enumerate}
    \item{Compute an intermediate velocity that does not necessarily
    satisfy the incompressibility constraint ($ \nabla \cdot \textbf u
    = 0$). This intermediate velocity is computed at each time step
    (ignoring the pressure term)}
    \item{Use the pressure $p$ to \textit{project} the intermediate
    velocity onto the space of divergence-free velocity fields.}
\end{enumerate}

Formally, the scheme reads

\begin{enumerate}
    \item Compute the intermediate velocity{$$\frac{\textbf u^*-\textbf
    u^n}{\Delta t} = - (\textbf u^n \cdot \nabla)\textbf u^n + \nu
    \nabla^2 \textbf u^n $$}
    \item Correct the velocity using the pressure {$$\frac{\textbf u^{n+1}
    -\textbf u^*}{\Delta t} = -\frac{1}{\rho} \nabla p^{n+1}$$} 
\end{enumerate}
There are things to note from this procedure. Firstly, the pressure $p$ is
time-invariant. This means that pressure $p$ is not affected by our
timestep $\Delta t$. Secondly, we can summarize that this procedure is a
kind of operator splitting.

This velocity vector, however, is not complying with our constraint of
zero divergence. Additionally, the solution of the second equation of this
scheme requires knowledge about $p$. We therefore rewrite the second
equation as
$$\textbf u^{n+1}=\textbf u^* - \frac{\Delta t}{\rho}\nabla p^{n+1}$$

Ensuring the absence of divergence $(\nabla \textbf u^{n+1}=0)$ allows us
to further extend this equation:
\begin{align}
\textbf u^{n+1} & =\textbf u^* - \frac{\Delta t}{\rho}\nabla p^{n+1} \\
\text{respecting divergence} \quad 0 &= \nabla \textbf u^* -
\frac{\Delta t}{\rho}\nabla^2 p^{n+1} \\
\Rightarrow \Delta p^{n+1}&= \frac{\rho}{\Delta t}(\nabla \textbf u^*)
\end{align}
The latter result is the Poisson equation for pressure.

We can see that the method is called a \textit{projection} because we
decompose $\textbf u^*$ into a divergence-free (solenoidal) part $\textbf
u_{sol}$ and irraotational part $\textbf u_{irrot}$. With the projection
step, $\textbf u_{irrot}=\nabla \phi$ is removed to ensure the remaining
part is negative divergence-free.

$$\textbf u = \textbf u_{\text{sol}}+\textbf u_{\text{irrot}}= 
\textbf u_{\text{sol}}+\nabla \phi $$

Since $\nabla \textbf u_{\text{sol}}=0$, the expression simplifies when w
e calculate the divergence of the total expression:
$$\Delta \textbf u = \Delta \phi $$

Here, $\Delta = \nabla^2$ represents the Laplacian operator. Consequently,
the solenoidal part can be calculated as  $\textbf u_{\text{sol}}
=\textbf u - \nabla \phi$.

This projection method is called \textit{Helmholtz-Hodge decomposition}.

We once again consider the Poisson equation:
$$\Delta p = \frac{\rho}{\Delta t}(\nabla \textbf u ^*) 
\quad \text{in the domain } \Omega =[0,1]^2$$

Assuming Neumann boundary conditions, 
$$\frac{\partial p}{\partial \textbf n}=0 \Leftrightarrow \nabla p \cdot 
\textbf n = 0 \quad \text{on }\partial \Omega$$

With these boundary conditions, we derive a solution that is not unique,
as $p$ is defined up to an additive constant. To solve the Poisson
equation, we equidistantly discretize the spatial domain.
$$y_i = i \Delta y \quad i=0,...,N_{y}-1$$
$$ z_j = j \Delta z \quad j=0,...,N_z -1$$

The Laplacian operator is approximated using the second-order finite
difference stencil
$$\Delta p = \nabla^2 p = \frac{\partial ^2 p}{\partial y^2}+
\frac{\partial^2 p}{\partial z^2} \approx 
\frac{p_{i+1,j}-2p_{i,j}+p_{i-1,j}}{\Delta y^2}+ 
\frac{p_{i,j+1}-2p_{i,j}+p_{i,j-1}}{\Delta z^2}$$

\subsection{Discretized Poisson Equation}

We now substitute the discretized Laplace operator into our Poisson
equation.
$$ \Delta p^{n+1}= \frac{rho}{\Delta t}(\nabla \textbf u^*) $$
Using the finite difference approximation of the Laplacian yields
$$-2(\frac{1}{\Delta y^2}+\frac{1}{\Delta z^2})p_{i,j}
+\frac{p_{i+1,j}+p_{i-1,j}}{\Delta y^2}+\frac{p_{i,j+1}+p_{i,j-1}}
{\Delta z^2}=\frac{\rho}{\Delta t}(\nabla \textbf u^*)_{i,j}$$

\todo{this is FFT now?}
To ensure that our numerical solution respects the boundary conditions, we
approximate derivatives using a central finite difference scheme
$$ f'_i \approx \frac{f_{i+1}-f_{i-1}}{2 \Delta x} $$
Each boundary of the domain then corresponds to a numerical constraint,
summarized in Table~\ref{tab:boundaryconditions}. 

\begin{table}\label{tab:boundaryconditions}
\begin{center}
\begin{tabular}{l c c c}
    \hline
    \textbf{Boundary} & \textbf{BC} & \textbf{numerical BC} & 
    \textbf{Resulting condition}\\
    \hline
    \hline 
    $\text{Bottom} (z=0,j=0)$ & $\frac{\partial p}{\partial z}=0$ 
    & $\frac{p_{i,1}-p_{i,-1}}{2\Delta z}=0$ & $p_{i,1}=p_{i,-1}$ \\
    $\text{Top} (z=1,j=N_z-1) $& $\frac{\partial p}{\partial z}=0$
    & $\frac{p_{i,N_z}-p_{i,-2}}{2\Delta z}=0$ & $p_{i,N_z}=p_{i,N_z-2}$ \\
    $\text{Left} (y=0,i=0)$ & $\frac{\partial p}{\partial y}=0$
    & $\frac{p_{1,j}-p_{-1,j}}{2\Delta z}=0$ & $p_{1,j}=p_{-1,j}$ \\
    $\text{Right} (y=1,i=N_y-1) $ & $\frac{\partial p}{\partial y}=0$
    & $\frac{p_{N_y,j}-p_{N_y-2,j}}{2\Delta z}=0$ 
    & $p_{N_y,j}=p_{N_y-2,j}$ \\
\end{tabular}
\caption{\textbf{Dsicretized Boundary Conditions (BC)} Boundary conditions
for the four boundaries of the two-dimensional domain $\Omega = [0,1]^2$.
Each row shows the analytical boundary condition, its numerical
approximation, and the resulting condition for the discrete pressure
field.}
\end{center}
\end{table}

\subsection{Numerical Procedure}
To summarize, the numerical procedure is given by the following steps:
\begin{enumerate}
    
    \item{Solve $\frac{\textbf u^*-\textbf u^n}{\Delta t} 
    = -(\textbf u^n\cdot \nabla)\cdot \textbf u^n + \nu \Delta
    \textbf n^n$} using a suitable integrator scheme 
    \item{Compute pressure $p$ by solving the discretized Poisson
    equation}
    \item{Solve $\frac{\textbf u^{n+1}-\textbf u^*}{\Delta t}
    =-\frac{1}{\rho}\nabla p$ and obtain $\textbf u^{n+1}$}
    
\end{enumerate}
\section{Rayleigh--Bénard convection}
So far we have neglected any forcing in the velocity equation of the
Navier-Stokes system. In this section, we introduce thermal forcing in
order to reproduce the Rayleigh--Bénard convection cells. Rayleigh--Bénard
convection occurs in many physical systems. For example, it appears in
climatology on a variety of spatial scales, and convective phenomena play
an important role in the dynamics of the Earth's atmosphere. In this
section we determine suitable initial conditions and parameters to
reproduce such convective behavior using our numerical solution.

\subsubsection{Physical Parameters}
To begin, we again consider the momentum equation. For our problem we
apply the \textit{Boussinesq approximation}. \todo{explain Boussinesque-
approximation} With this approximation, the Navier--Stokes equations
become
$$\partial_t \textbf{u} + (\textbf{u})\cdot \nabla 
+ \frac{1}{\rho}\nabla p = \nu \nabla \textbf u + T \beta e_z $$

As we are interested in a suitable set of parameters, we briefly discuss
the physical meaning of the most important ones. Consider the viscosity
$\nu$ for example. If $\nu$ is too large (corresponding to strong
damping), the fluid remains in a purely conductive state and no convective
cells develop. The buoyancy parameter $\beta$ is also crucial. If $\beta$
is too small, heat will simply diffuse through the fluid instead of
generating motion due to buoyancy forcing. Thus, these parameters are
essential for a successful recreation of the convective cells. While
viscosity $\nu$ determines how strongly viscous effects damp the motion of
the fluid, the parameter $\beta$ controls the strength of the coupling
between temperature and velocity.

In a dimensionless formulation, the onset of convection is determined by
the \textit{Rayleigh number}. In our simplified setting, the Rayleigh
number is directly proportional to $\frac{\beta}{\nu}$. Therefore, $\beta$
must be sufficiently large enough (or $\nu$ sufficiently small) for
buoyancy forces to overcome viscous damping and allow the formation of
convection cells.

\subsection{Initial conditions}
To trigger the formation of convection cells, it is not sufficient to
choose appropriate physical parameters. Suitable initial conditions are
also required. As symmetry in the initial conditions could hinder the
formation of convective cells, we cannot start with a perfectly uniform or
stratified fluid. Therefore, we introduce small perturbations into the
initial conditions.

Firstly, we initialize the temperature field with a conductive profile.
That is, the initial temperature distribution $T_0(y,z)$ approximately
follows the linear profile $1-z.$ This corresponds to a fluid layer heated
from below and cooled from above.

Secondly, we want to break the symmetry in the velocity field by
introducing small perturbations. We chose random noise that is used to
initialize the velocity field $\textbf u_0(y,z)$. These perturbations will
be amplified by the the buoyancy term $T \beta e_z$ and, if the Rayleigh
number is high enough, eventually lead to the emergence of Rayleigh
-Bernard convection.

Formally, the described set of initial conditions can be described in the
following way. On our domain $\Omega=[0,1]^2,$ we introduce time-invariant
temperatures (constant heating across time): In a perfectly still fluid
($\textbf u=0$) the temperature settles into a linear conductivity
profile.
$$T(x,y)= 1-z$$
This set up implies $T=1$ at $z=0$ (bottom) and $T=0$ at $z=1$ (top).
Note that this state satisfies the equilibrium condition meaning 
$$
\partial _t T + \textbf u \Delta T = \Delta T,
$$
where $\Delta T=0$ and $\textbf u \cdot \Delta t=0$.

To introduce the mentioned perturbations is the temperature field, we add
noise to these initial conditions.
$$T=(1-z)+ \epsilon $$

To understand if such perturbations will grow or vanish, we substitute the
temperature field into the momentum equations and linearize them. This is
performed by dropping nonlinear terms like $\nu \cdot \nabla \epsilon$ or
$(\textbf u \cdot \nabla) \textbf u$

That is how we obtain the temperature equations
\begin{align}
\partial_t + u \cdot \nabla (1-z)&=\Delta \epsilon \\
\nabla (1-z)&=-e_z \\
\end{align}

This leads to $\textbf u \cdot (1-z) = -w$
In total, we obtain the linearized heat equation
$$\partial_t \epsilon = w + \Delta \epsilon$$

This relation shows that upward motion $w>0$ acts as the source term that
increases the local temperature perturbation.

If we neglect nonlinear terms in the momentum equation and decompose
pressure into a background component and perturbation $p'$, the vertical
component of the momentum equation becomes:
$$\partial_t w + \frac{1}{\rho}\partial_z p'
= \nu \Delta w + \epsilon \beta$$

Here, the term $\epsilon \beta $ represents the buoyancy force. It drives
vertical motion $w$, which in turn increases the temperature perturbation
$\epsilon$. Thus, stability of the system is determined by the competition
between buoyancy and dissipative effects.

The perturbations are influenced by two competing mechanisms:
\begin{enumerate}
    \item{\textbf{Driving forces}: The buoyancy term ($T\beta e_z$), which
    amplifies the disturbances as warm fluid rises.}
    \item{\textbf{Damping forces}: The viscosity $\nu \Delta \textbf u$
    and thermal diffusion $\Delta T$, which smooth velocity and
    temperature gradients.}
\end{enumerate} 

In conclusion, we see that the growth rate of the perturbations $\epsilon$
become positive only when the buoyancy term overcomes these damping
effects. If $\beta $ is too small, diffusion is more important, and the
perturbation decays back to the conductive state.